\section{Verified Artifact Facts}
\label{app:facts}

Every figure below was measured against the repository, on a clean build,
rather than asserted.  Measurements are from an Apple M5 (10 cores, 32~GB)
with Mathlib's compiled artifacts supplied by \texttt{lake exe cache get}.

\begin{table}[h]
\centering
\small
\begin{tabular}{ll}
\toprule
Item & Value \\
\midrule
Lean toolchain & \texttt{leanprover/lean4:v4.26.0} \\
Mathlib version & \texttt{v4.26.0} \\
Mathlib revision & \texttt{2df2f0150c275ad53cb3c90f7c98ec15a56a1a67} \\
Lake libraries & 4 \\
Lean files & 30 \\
Lines of Lean & 6{,}499 \\
Lines excluding comments and blanks & 4{,}471 \\
Hand-written \texttt{theorem} declarations & 274 \\
Hand-written \texttt{def}/\texttt{abbrev} declarations & 90 \\
\texttt{lake build} & succeeds, 829 jobs \\
Clean build, Mathlib cached & 2\,min 2\,s \\
Slowest module (\texttt{FriedbergMuchnik/CE.lean}) & 106\,s \\
Declarations checked with \texttt{\#print axioms} & 37 \\
Axiom footprint, all 37 & \texttt{propext}, \texttt{Classical.choice}, \texttt{Quot.sound} \\
Axioms declared by the development & 0 \\
\texttt{sorry} / \texttt{admit} & 0 \\
\texttt{native\_decide} / \texttt{unsafe} / \texttt{partial def} & 0 \\
\texttt{noncomputable} declarations & 2 \\
Declarations with raised \texttt{maxHeartbeats} & 11 \\
Repository license & none \\
CI workflow & none \\
\bottomrule
\end{tabular}
\caption{Artifact facts, measured at the revision described in
Appendix~\ref{app:artifact}.}
\end{table}

\begin{table}[h]
\centering
\small
\begin{tabular}{lrrr}
\toprule
Library & Files & Lines & Code lines \\
\midrule
\texttt{OracleComputability} & 14 & 2{,}866 & 2{,}000 \\
\texttt{FriedbergMuchnik} & 8 & 2{,}002 & 1{,}494 \\
\texttt{SacksSplitting} & 7 & 1{,}548 & 948 \\
\texttt{PriorityArguments} & 1 & 83 & 29 \\
\bottomrule
\end{tabular}
\caption{Size by library.  ``Code lines'' excludes comments and blanks.}
\end{table}

The two \texttt{noncomputable} declarations are the classical
characteristic function \texttt{charFun}, which appears only in the
statements of the reducibility vocabulary, and \texttt{limUse}, which
appears only in the verification of the second theorem.  Neither occurs in
a construction.  The raised heartbeat budgets are elaboration limits in the
two computable-enumerability files and cannot affect the axiom footprint,
which is the ground truth that no logical shortcut was taken.

\paragraph{The constructions run.}
Both stage machines are executable Lean functions, and evaluating them is a
stronger check than the absence of \texttt{sorry}: it shows the definitions
are not vacuously satisfiable stubs.  Evaluating the Friedberg--Muchnik
machine --- the one the priority proof reasons about, not the shadow
implementation used for the computability proof --- at stage 20 gives
enumerations $\{12,8,6,4,0\}$ and $\{13,10,1\}$ for the two sets, a
freshness counter of $14$, and strategy records in which requirements 8 and
10 hold restraints 9 and 11 respectively.  Evaluating the splitting machine
on the code that halts everywhere routes the first eight elements into a
single half, all restraints remaining zero, since a reduction that always
answers zero never agrees with a set that is accumulating members.
